\section{Mathematical Content}\label{sec:mathematical-content}

\subsection{Inline Mathematics}\label{subsec:inline-mathematics}
LaTeX excels at typesetting mathematics. Inline math is enclosed in \texttt{\$...\$} or \texttt{\(...\)}:

\begin{itemize}
    \item The Pythagorean theorem: $a^2 + b^2 = c^2$
    \item Euler's identity: $e^{i\pi} + 1 = 0$
    \item A fraction in text: \(f(x) = \frac{x^2 - 1}{x - 1}\) simplifies to $f(x) = x + 1$ for $x \neq 1$
    \item Subscripts and superscripts: $x_i^2 + y_j^2 = z_{ij}^2$
\end{itemize}

\subsection{Display Equations}\label{subsec:display-equations}
For standalone equations, use \texttt{\[...\]} or the \texttt{equation} environment:

\begin{verbatim}
\[
    \int_{a}^{b} f(x) \, dx = F(b) - F(a)
\]
\end{verbatim}

Which produces:
\[
    \int_{a}^{b} f(x) \, dx = F(b) - F(a)
\]

\subsection{Numbered Equations}\label{subsec:numbered-equations}
The \texttt{equation} environment automatically numbers equations:

\begin{equation}
    E = mc^2
    \label{eq:einstein}
\end{equation}

You can reference equations using \texttt{\\ref} or \texttt{\\eqref}: Equation~\eqref{eq:einstein} shows Einstein's famous equation.

\subsection{Mathematical Symbols}\label{subsec:mathematical-symbols}
LaTeX provides an extensive set of mathematical symbols:

\subsubsection{Greek Letters}\label{subsubsec:greek-letters}
$\alpha, \beta, \gamma, \delta, \epsilon, \varepsilon, \zeta, \eta, \theta, \vartheta$

$\pi, \varpi, \rho, \varrho, \sigma, \varsigma, \tau, \upsilon, \phi, \varphi$

$\Gamma, \Delta, \Theta, \Lambda, \Xi, \Pi, \Sigma, \Upsilon, \Phi, \Psi, \Omega$

\subsubsection{Operators}\label{subsubsec:operators}
$+, -, \times, \div, \pm, \mp, \cdot, \star, \circ, \bullet$

\subsubsection{Relations}\label{subsubsec:relations}
$=, \neq, \approx, \sim, \simeq, \cong, \equiv, <, \leq, >, \geq, \ll, \gg$

\subsubsection{Set Symbols}\label{subsubsec:set-symbols}
$\in, \notin, \subset, \subseteq, \supset, \supseteq, \cap, \cup, \emptyset$

\subsubsection{Calculus}\label{subsubsec:calculus}
$\sum_{i=1}^{n}, \prod_{i=1}^{n}, \int_{a}^{b}, \oint, \nabla, \partial, \infty, \lim_{x \to \infty}$

\subsection{Multiple Equations}\label{subsec:multiple-equations}
The \texttt{align} environment from \texttt{amsmath} aligns multiple equations:

\begin{align}
    (a+b)^2 &= (a+b)(a+b) \\
    &= a^2 + ab + ba + b^2 \\
    &= a^2 + 2ab + b^2
    \label{eq:binomial-expansion}
\end{align}

\subsection{Matrix Notation}\label{subsec:matrix-notation}
Various matrix environments are available:

\subsubsection{Basic Matrix}\label{subsubsec:basic-matrix}
\[
\begin{matrix}
a & b & c \\
d & e & f \\
g & h & i
\end{matrix}
\]

\subsubsection{Parenthesized Matrix}\label{subsubsec:parenthesized-matrix}
\[
\begin{pmatrix}
a & b & c \\
d & e & f \\
g & h & i
\end{pmatrix}
\]

\subsubsection{Determinant}\label{subsubsec:determinant}
\[
\begin{vmatrix}
a & b & c \\
d & e & f \\
g & h & i
\end{vmatrix}
\]

\subsection{Piecewise Functions}\label{subsec:piecewise-functions}
The \texttt{cases} environment handles piecewise definitions:

\[
f(x) = \begin{cases}
    x^2 & \text{if $x \geq 0$} \\
    -x^2 & \text{if $x < 0$}
\end{cases}
\]

\subsection{Advanced Equation Formatting}\label{subsec:advanced-equation-formatting}
For complex equations, additional environments help:

\subsubsection{Fractions and Binomials}\label{subsubsec:fractions-binomials}
\[
\frac{1}{1 + \frac{1}{1 + \frac{1}{1 + x}}} \quad \text{and} \quad
\binom{n}{k} = \frac{n!}{k!(n-k)!}
\]

\subsubsection{Multi-line Equations}\label{subsubsec:multi-line-equations}
Use \texttt{aligned} for equations that need alignment within a larger equation:

\[
\begin{aligned}
S &= \sum_{i=1}^{n} i^2 \\
&= \frac{n(n+1)(2n+1)}{6}
\end{aligned}
\]

\subsubsection{Equation Arrays}\label{subsubsec:equation-arrays}
The \texttt{array} environment offers complete flexibility:

\[
\left\{
\begin{array}{rcl}
x + y + z &=& 10 \\
xy + yz + zx &=& 24 \\
xyz &=& 40
\end{array}
\right.
\]

\subsection{Theorem Environments}\label{subsec:theorem-environments}
Mathematical documents often include theorems and lemmas:

\begin{verbatim}
% Add this to your preamble:
% \usepackage{amsthm}
% \newtheorem{theorem}{Theorem}[section]
% \newtheorem{lemma}[theorem]{Lemma}
\end{verbatim}

A theorem would then be written as:

\begin{verbatim}
\begin{theorem}
For any right triangle with sides $a$, $b$ and hypotenuse $c$, $a^2 + b^2 = c^2$.
\end{theorem}
\end{verbatim}

% Note: Proof environment requires amsthm package
% To use the proof environment, add \usepackage{amsthm} to your preamble